\documentclass[11pt,a4paper]{article}
\usepackage[T1]{fontenc}
\usepackage[utf8]{inputenc}
\usepackage{amsmath,amssymb,xcolor}
\usepackage[margin=25mm]{geometry}
\usepackage[colorlinks=true,urlcolor=blue]{hyperref}
\definecolor{nexorablue}{HTML}{2459B5}
\setlength{\parindent}{0pt}
\setlength{\parskip}{7pt}
\setlength{\emergencystretch}{2em}
\renewcommand{\familydefault}{\sfdefault}
\begin{document}
\makeatletter
\def\@oddfoot{\footnotesize Download for free at \url{https://openstax.org/details/books/university-physics-volume-2}.\hfil\thepage}
\let\@evenfoot\@oddfoot
\makeatother
{\small\bfseries\color{nexorablue} NEXORA CPGE}\par
{\small MPSI / OST-SUP-14-04 / Énoncé}\par
{\LARGE\bfseries Établissement du courant dans une bobine}\par
\vspace{0.25em}\hrule\vspace{0.6em}

Un générateur de $2{,}0\,\mathrm V$ est appliqué à $t=0$ à une boucle série comportant une résistance de $4{,}0\,\mathrm\Omega$ et une inductance de $4{,}0\,\mathrm H$. Le courant initial est nul. Déterminer la constante de temps, puis le courant et la valeur absolue de la tension inductive à $t=0^+$, $t=2\tau$ et au régime permanent.

\bigskip
{\footnotesize Adaptation française par Nexora CPGE. Source : Samuel J. Ling, Jeff Sanny et William Moebs (OpenStax, Rice University),
\href{https://d3bxy9euw4e147.cloudfront.net/oscms-prod/media/documents/UniversityPhysicsVolume2-OP.pdf\#page=652}{\textit{University Physics Volume 2, édition 2018}}, § 14.4, exemple 14.4, p. 642–643.
\href{https://creativecommons.org/licenses/by/4.0/}{CC BY 4.0}.}
\end{document}
